(1) 以\(x=\Delta/(k_BT)\)计，单粒子配分函数\(z=g_0+g_1e^{-x}\)，全体系\(Z_N=z^N\)。定义\(r=(g_1/g_0)e^{-x}\)，激发态概率\(p_1=g_1e^{-x}/z=r/(1+r)\)，而基态概率\(1-p_1\)。简并度乘入Boltzmann权重，不能误把每一级当成一个微观态。
(2) 热内能\(U=-\partial_\beta\ln Z_N=N\Delta p_1\)，又\(\partial_Tp_1=p_1(1-p_1)\Delta/(k_BT^2)\)，因此\(C_V=Nk_Bx^2r/(1+r)^2\)。等价写成\(C_V=N\operatorname{Var}(E)/(k_BT^2)\)，因为\(\operatorname{Var}(E)=\Delta^2p_1(1-p_1)\)；这给出由正则涨落得到的第二条检查路线。
(3) 若\(g_0=g_1\)，有\(r=e^{-x}\)，\(C_V/(Nk_B)=x^2/[4\cosh^2(x/2)]\)。对\(x>0\)求导，驻点条件\(2/x-\tanh(x/2)=0\)，即\(x\tanh(x/2)=2\)。左边从0严格递增至无穷，故有唯一正根\(x_*\simeq2.399357\)，峰温\(T_*=\Delta/(k_Bx_*)\)；该热容峰是有限能级间热激发，不属于相变潜热。
(4) 低温\(x\gg1\)时\(r\ll1\)，\(C_V\sim Nk_B(g_1/g_0)x^2e^{-x}\)，指数冻结；高温\(x\ll1\)时\(r\to g_1/g_0\)，\(C_V\sim Nk_Bx^2g_0g_1/(g_0+g_1)^2\to0\)。简并度改变峰的高度和位置，但不会在有限能级的独立粒子模型中产生热力学奇异性。
